%%
%% CSD Assignment 1 - Final Technical Report
%% Based on the ACM "acmart" manuscript (single-column) style.
%% Course: Computer Science and Design (CSD), IIT Bhilai
%%
\documentclass[manuscript,screen,review]{acmart}

\AtBeginDocument{%
  \providecommand\BibTeX{{Bib\TeX}}}

\usepackage{booktabs}
\usepackage{amsmath}
\usepackage{graphicx}
\usepackage{listings}
\usepackage{xcolor}

\begin{document}

%% =================== TITLE & AUTHOR BLOCK =========================
\title{AquaPlan AI: AI-assisted Village Pond Planning, Catchment Delineation, and Runoff Harvesting System}
\subtitle{CSD Assignment 1 -- Final Technical Report}

\author{Gowrabathuni Nikhilesh}
\affiliation{%
  \institution{Department of Computer Science and Engineering, IIT Bhilai}
  \city{Bhilai}
  \state{Chhattisgarh}
  \country{India}}
\email{gowrabathunin@iitbhilai.ac.in}

\renewcommand{\shortauthors}{G. Nikhilesh}

%% =================== ABSTRACT ======================================
\begin{abstract}
Rural water conservation in semi-arid and agrarian regions of India heavily depends on decentralized rainwater harvesting structures such as village ponds. However, conventional manual site selection relies on coarse topographic heuristics and rudimentary site inspections, frequently resulting in inadequate catchment yield, severe embankment breach risks, or rapid siltation. In this paper, we present \textbf{AquaPlan AI}, an automated, end-to-end geospatial and hydrological planning platform designed to identify optimal village pond locations, delineate topographical catchment drainage basins, and estimate collectible rainwater runoff volume. The system integrates digital contour extraction from high-resolution Keyhole Markup Language (KML) datasets, spatial projection into Universal Transverse Mercator (UTM) coordinates, continuous Digital Elevation Model (DEM) surface interpolation, and Deterministic 8-Neighbor (D8) hydrological flow accumulation algorithms. candidate sinks are systematically ranked through a multi-criteria objective function balancing contributing watershed area, local terrain gradient, upland elevation advantage, and river buffer avoidance. Furthermore, the system incorporates empirical runoff modeling using 30-year rainfall normals from the India Meteorological Department (IMD) for the Bhilai/Durg district (mean annual precipitation: 1220~mm) and Central Ground Water Board (CGWB) soil runoff coefficients. The backend is deployed on an asynchronous FastAPI microservice architecture with self-healing process supervision, while the frontend delivers an interactive, responsive GIS web dashboard utilizing Leaflet for custom bounding-box land area selection and dynamic vector overlays. Experimental evaluations across 2,710 contour features encompassing the IIT Bhilai study area demonstrate sub-50~ms cached query response times and 2.54~s processing times for dynamic sub-watershed queries, providing a robust, fault-tolerant tool for rural water engineering.
\end{abstract}

\keywords{Village Pond Planning, Digital Elevation Model (DEM), D8 Flow Routing, Catchment Delineation, Rainwater Harvesting, Runoff Volume Estimation, Web GIS, FastAPI}

\maketitle

%% =====================================================================
\section{Introduction}
\label{sec:intro}
Water scarcity remains one of the most pressing socio-economic challenges across rural India, particularly in agrarian states such as Chhattisgarh where agriculture accounts for over 70\% of local livelihood. In rainfed agro-ecosystems, erratic monsoon distributions lead to severe seasonal droughts interspersed with intense flash runoff. Decentralized rainwater harvesting through earthen village ponds (\textit{Pokhar} or \textit{Talab}) offers an ecologically sustainable solution to recharge groundwater aquifers, provide supplementary irrigation during dry spells, and support livestock. However, the efficacy of a rainwater harvesting pond depends critically on its spatial placement. A poorly situated pond with insufficient contributing catchment fails to fill, while a pond placed in high-energy river thalwegs is vulnerable to catastrophic embankment failure and rapid siltation. 

The primary objective of this project is to develop an automated, AI-assisted geospatial web platform that processes topographical contour maps, computes digital elevation models, delineates contributing drainage basins via hydrological flow routing, and recommends optimal pond locations along with expected water yields. The rest of this report is organized as follows: Section~\ref{sec:requirements} defines functional and non-functional requirements. Section~\ref{sec:architecture} describes the system architecture and software stack. Section~\ref{sec:methodology} explains the algorithmic methodology of terrain analysis and runoff estimation. Section~\ref{sec:implementation} outlines backend, frontend, and storage implementations. Section~\ref{sec:csd-themes} maps system modules to core Computer System Design (CSD) principles. Section~\ref{sec:results} presents quantitative evaluation results, followed by discussions in Section~\ref{sec:discussion}, AI declarations in Section~\ref{sec:ai-usage}, and deployment URLs in Appendix~\ref{sec:appendix}.

\subsection{Motivation}
Manual site selection conducted by village administrators or local panchayat engineers typically suffers from three fundamental bottlenecks:
\begin{enumerate}
    \item \textbf{Topographical Complexity:} Visual inspection cannot accurately identify subtle depressions, micro-watersheds, and overland flow paths across low-relief agricultural plains.
    \item \textbf{Decoupled Meteorological Data:} Field planners rarely integrate localized, multi-decadal historical rainfall records and soil infiltration parameters to determine whether a catchment can hydrologically sustain the planned excavation volume.
    \item \textbf{Lack of Interactive Computational Tools:} Traditional GIS software suites (e.g., ArcGIS, QGIS) require high computational overhead, proprietary licenses, and specialized training, making them inaccessible for real-time field deliberation. An automated, web-accessible GIS platform eliminates these barriers by providing instantaneous hydrological insights upon selecting an arbitrary parcel of land on an interactive map.
\end{enumerate}

\subsection{Scope of the Project}
The system specifically encompasses: (i) parsing and spatial transformation of contour data from KML/KMZ formats; (ii) continuous DEM generation via linear and nearest-neighbor spatial interpolation; (iii) automated D8 flow direction and accumulation routing; (iv) multi-criteria ranking of candidate pond sites excluding active river channels; (v) convex-hull boundary delineation of contributing catchment basins; (vi) empirical runoff volume calculation based on IMD rainfall normals and soil classifications; and (vii) an interactive web GIS dashboard supporting dynamic bounding-box land parcel selection. The project does \emph{not} cover geotechnical soil shear-strength testing, structural reinforced concrete spillway design, or cadastral land-ownership legal verification.

%% =====================================================================
\section{Problem Statement and Requirements}
\label{sec:requirements}
The assignment mandates developing an operational web application that empowers users to visualize terrain contours, interactively select a land parcel on a map, and instantaneously generate optimal pond placement recommendations, delineated catchment boundaries, and expected collectible water volumes. Table~\ref{tab:requirements} maps each functional requirement to its corresponding architectural component.

\begin{table}[h]
  \caption{Functional requirements and system module mapping}
  \label{tab:requirements}
  \begin{tabular}{p{4.2cm}p{4.5cm}p{4.5cm}}
    \toprule
    Requirement & Implemented in Module & Technical Mechanism \\
    \midrule
    Satellite imagery display & \texttt{static/app.js}, \texttt{index.html} & Esri World Imagery \& OSM tile layers in Leaflet \\
    Contour map visualization & \texttt{generate\_assets.py}, \texttt{static/app.js} & GeoJSON vector layer with elevation color ramp \\
    Interactive land area selection & \texttt{static/app.js} (Bounding Box Tool) & Leaflet drag-and-drop bounding rectangle \& preset zones \\
    Catchment area estimation & \texttt{terrain\_analysis.py} & D8 flow routing \& upstream pixel cell accumulation \\
    Historical rainfall query & \texttt{main.py} (\texttt{/api/rainfall}) & IMD 30-year precipitation normals (1220~mm annual) \\
    Runoff volume estimation & \texttt{terrain\_analysis.py}, \texttt{main.py} & Rational runoff formulation ($V = C \cdot P \cdot A$) \\
    Pond depth/storage recommendation & \texttt{terrain\_analysis.py} & Hydraulic sizing (3.5~m depth, 1.5:1 side slopes) \\
    Combined overlay on map & \texttt{static/app.js}, \texttt{static/style.css} & Pulsing beacon marker, translucent polygon, and badges \\
    \bottomrule
  \end{tabular}
\end{table}

\subsection{Non-Functional Requirements}
\begin{itemize}
    \item \textbf{Performance and Latency:} Precomputed baseline study area queries must complete in under 100~ms; dynamic bounding-box queries on subregions must execute within 3~seconds.
    \item \textbf{Availability and Fault Tolerance:} The backend services must run continuously in the background under process supervision, automatically restarting upon fatal runtime anomalies.
    \item \textbf{Port Compliance and Network Accessibility:} The application must adhere strictly to assigned system port constraints (ports 2000--7000), exposing the full web interface across the internal institutional network at assigned port mappings.
    \item \textbf{Responsiveness and Usability:} The web frontend must operate smoothly across varying desktop screen resolutions, utilizing asynchronous REST communication to prevent UI freezing during intensive geospatial matrix computations.
\end{itemize}

%% =====================================================================
\section{System Architecture and High-Level Design}
\label{sec:architecture}
The AquaPlan AI platform follows a decoupled, three-tier microservice architecture comprising the Presentation Layer (Web GIS Client), Application Layer (FastAPI REST Engine), and Data Processing/Hydrology Layer (NumPy, SciPy, PyProj). Figure~\ref{fig:architecture} illustrates the end-to-end data pipeline.

\begin{figure}[h]
  \centering
  \fbox{
  \begin{minipage}{0.9\linewidth}
  \centering
  \small
  \textbf{[Interactive Web GIS Frontend (Leaflet / Vanilla JS)]}\\
  $\downarrow$ \textit{User selects bounding box on map / uploads KML} (HTTP POST)\\
  \textbf{[FastAPI Asynchronous Gateway (Uvicorn / Ports 3000 \& 6000)]}\\
  $\downarrow$ \textit{Request dispatched to Hydrological Processing Engine}\\
  \textbf{[Terrain Analysis Module (\texttt{terrain\_analysis.py})]}\\
  1. KML XML Parsing (\texttt{lxml.etree}) $\rightarrow$ Contours $(x, y, z)$\\
  2. Coordinate Transformation (\texttt{pyproj}) $\rightarrow$ WGS84 to UTM Zone 44N\\
  3. Continuous DEM Generation (\texttt{scipy.interpolate.griddata})\\
  4. Hydrological Flow Routing (D8 Matrix Operations)\\
  5. Sinks Delineation \& River Avoidance Filtering\\
  6. Watershed Convex Hull Extraction (\texttt{scipy.spatial.ConvexHull})\\
  7. Runoff Hydrology \& Storage Sizing ($V = C \cdot P \cdot A$)\\
  $\downarrow$ \textit{JSON Response with GeoJSON Catchment Polygon \& Recommendations}\\
  \textbf{[Client-side Overlay Visualization \& Dynamic Dashboard Update]}
  \end{minipage}
  }
  \caption{High-level architectural pipeline of AquaPlan AI.}
  \label{fig:architecture}
  \Description{Block diagram showing data flow from the web GIS client to FastAPI, terrain analysis engine, and back.}
\end{figure}

\subsection{Technology Stack}
The system implementation employs:
\begin{itemize}
    \item \textbf{Backend Framework:} \textbf{FastAPI} (v0.141.1) running on \textbf{Uvicorn} ASGI web server for high-throughput, non-blocking asynchronous request dispatching.
    \item \textbf{Scientific \& Geospatial Libraries:} \textbf{NumPy} (v2.5.2) for vectorized raster manipulation; \textbf{SciPy} for multidimensional Delaunay grid interpolation (\texttt{griddata}) and convex hull delineation (\texttt{ConvexHull}); \textbf{PyProj} for geodesic coordinate transformations between EPSG:4326 (WGS84) and EPSG:32644 (UTM Zone 44N); and \textbf{lxml} for high-speed XML contour parsing.
    \item \textbf{Frontend GIS \& UI:} Modern Vanilla JavaScript (ES6+), HTML5 semantic structure, Vanilla CSS with custom glassmorphic design tokens, and \textbf{Leaflet.js} (v1.9.4) for interactive mapping, tile rendering, and vector polygon overlays.
    \item \textbf{Process Supervision:} Custom Python watchdog supervisor (\texttt{start\_daemon.py}) providing daemonized process monitoring, automatic fault recovery, and dual-port load serving.
\end{itemize}

%% =====================================================================
\section{Methodology}
\label{sec:methodology}

\subsection{Terrain and Elevation Analysis}
Topographic elevation information is sourced from Keyhole Markup Language (KML) vector datasets comprising discrete contour isolines. The parsing module utilizes \texttt{lxml.etree} to extract all \texttt{<Placemark>} elements containing line coordinates and elevation tags. Because longitude/latitude degrees represent non-uniform spatial distances, raw WGS84 coordinates $(\lambda, \phi)$ are projected into metric Universal Transverse Mercator (UTM Zone 44N) planar coordinates $(x, y)$ using the standard ellipsoidal transformation:
\begin{equation}
    (x, y) = \mathcal{T}_{\text{WGS84} \rightarrow \text{UTM}}(\lambda, \phi)
\end{equation}
A regular 2D grid matrix $\mathbf{G}_{x, y}$ is initialized at a uniform grid resolution of $\Delta x = \Delta y = 2.0\text{ meters}$. To transform discrete contour polyline vertices into a continuous Digital Elevation Model (DEM), spatial 2D surface interpolation is performed using linear barycentric interpolation on the Delaunay triangulation of the input coordinates. Residual boundary grid cells outside the convex hull of input contours are seamlessly imputed using nearest-neighbor Euclidean interpolation:
\begin{equation}
    Z(x_0, y_0) = \sum_{i=1}^3 w_i Z_i, \quad \text{where } \sum w_i = 1
\end{equation}

\subsection{Catchment Area Delineation via D8 Flow Routing}
Overland hydrological drainage pathways are modeled using the standard \textbf{Deterministic 8-Neighbor (D8)} flow routing algorithm. For every raster cell $(r, c)$ in the DEM matrix $\mathbf{Z}$, the algorithm evaluates the steepest downward slope among its eight immediate orthogonal and diagonal neighbors:
\begin{equation}
    \Delta Z_{d} = \frac{Z(r, c) - Z(r + \Delta r_d, c + \Delta c_d)}{d_d}, \quad d \in \{1, \dots, 8\}
\end{equation}
where $d_d = \Delta x$ for orthogonal steps and $d_d = \sqrt{2}\Delta x$ for diagonal steps. Each cell routes its flow to the neighbor exhibiting maximum positive $\Delta Z_d$. Cells with no downward neighbors are identified as topographical depression sinks. An iterative topological path resolution with cycle detection maps each raster cell $i$ to its terminal sink $S(i)$.

To prevent siting ponds inside active, high-discharge river channels or seasonal streams, a flow accumulation raster $\mathbf{A}_{\text{flow}}$ is computed by sorting all elevation cells in descending order and accumulating upstream drainage cell counts:
\begin{equation}
    \mathbf{A}_{\text{flow}}[S(i)] = \mathbf{A}_{\text{flow}}[S(i)] + \mathbf{A}_{\text{flow}}[i]
\end{equation}
Cells exceeding the 97th percentile of flow accumulation ($\mathbf{A}_{\text{flow}} > \tau_{\text{river}}$) are classified as river network paths and strictly penalized. Candidate pond locations are selected by ranking depression sinks using a multi-criteria objective score:
\begin{equation}
    \text{Score} = 0.3 \cdot \bar{Z}_{\text{norm}} + 0.4 \cdot \left[ \frac{A_k}{A_{\max}} \left(1 - \frac{A_k}{A_{\max}}\right) \right] + 0.3 \cdot \left( 1 - \frac{\mathbf{A}_{\text{flow}}}{\tau_{\text{river}}} \right)
\end{equation}
The contributing catchment polygon is then delineated by extracting all cells draining into the optimal sink $S^*$ and computing the two-dimensional planar boundary polygon using the Quickhull algorithm (\texttt{ConvexHull}).

\subsection{Rainfall Data Integration and Runoff Volume Estimation}
Historical rainfall data is integrated from the India Meteorological Department (IMD) 30-year climatological normals for the Bhilai/Durg district station (Lat: 21.2167°N, Lon: 81.3833°E). The long-term mean annual precipitation is $P_{\text{annual}} = 1220.0\text{ mm}$, with $85.6\%$ ($1044.3\text{ mm}$) occurring during the southwest monsoon season (June to September). 

Collectible surface runoff volume $V_{\text{runoff}}$ contributing to the pond is quantified using the standard Rational Method endorsed by the Central Ground Water Board (CGWB) and Indian Standard Code IS 14961:
\begin{equation}
    V_{\text{annual}} = C \cdot \left( \frac{P_{\text{annual}}}{1000} \right) \cdot A_{\text{catchment}} \quad [\text{m}^3]
\end{equation}
where $A_{\text{catchment}}$ is the contributing catchment area in $\text{m}^2$, and $C$ is the dimensionless runoff coefficient governed by surface soil permeability. For the prevailing black cotton and clayey soils of the Chhattisgarh plains, $C = 0.38\text{--}0.45$ (with loam at $0.35$ and sandy loam at $0.25$).

Hydraulic sizing recommendations for the pond embankment specify:
\begin{equation}
    V_{\text{design}} = \min(0.75 \cdot V_{\text{monsoon}}, 60000\text{ m}^3), \quad A_{\text{surface}} = \frac{V_{\text{design}}}{d_{\text{effective}}}
\end{equation}
with an optimal excavation depth of $d = 3.5\text{ m}$ (comprising $0.5\text{ m}$ freeboard, $2.5\text{ m}$ active storage, and $0.5\text{ m}$ silt dead storage) and side slopes of $1.5:1$ (H:V).

%% =====================================================================
\section{Implementation}
\label{sec:implementation}

\subsection{Backend and API Design}
The backend service is implemented in Python 3.12 utilizing FastAPI and Uvicorn. All endpoints conform strictly to RESTful design conventions and return structured JSON schemas. Table~\ref{tab:api-endpoints} details the exposed REST interface.

\begin{table}[h]
  \caption{AquaPlan AI REST API Endpoints}
  \label{tab:api-endpoints}
  \begin{tabular}{p{2.6cm}p{4.2cm}p{6.5cm}}
    \toprule
    HTTP Method & Endpoint Path & Description and Payload Contract \\
    \midrule
    \texttt{GET} & \texttt{/} & Serves the interactive Leaflet Web GIS client application \\
    \texttt{GET} & \texttt{/api/health} & Service uptime, active ports (3000, 6000), host status \\
    \texttt{GET} & \texttt{/api/info} & Geospatial bounding box, elevation range, and contour summary \\
    \texttt{GET} & \texttt{/api/contours} & Returns pre-simplified GeoJSON vector contour isolines \\
    \texttt{GET} & \texttt{/api/baseline} & Instant cached baseline analysis of full IIT Bhilai study area \\
    \texttt{GET} & \texttt{/api/rainfall} & Monthly IMD rainfall distribution and soil runoff coefficients \\
    \texttt{POST} & \texttt{/api/analyze-bounds} & Analyzes custom bounding box: \texttt{\{min\_lat, max\_lat, min\_lon, max\_lon, runoff\_coef\}} \\
    \texttt{POST} & \texttt{/api/calculate-water-volume} & Dynamic recalculation of yield and pond sizing on parameter change \\
    \texttt{POST} & \texttt{/analyzeContour} & Multipart file upload of custom \texttt{.kml} or \texttt{.kmz} contour map \\
    \bottomrule
  \end{tabular}
\end{table}

\subsection{Frontend and Visualization}
The web interface features an ultra-modern dark glassmorphism theme developed with modular CSS variables and vanilla JavaScript. Key frontend capabilities include:
\begin{itemize}
    \item \textbf{Dual Base Mapping:} Seamless switching between high-resolution satellite imagery (Esri World Imagery) and topographic street basemaps (OpenStreetMap).
    \item \textbf{Contour Isolines Overlay:} Dynamic vector lines color-coded along a smooth spectral gradient from cyan (267~m) through green and amber to red (298~m), with interactive elevation hover tooltips.
    \item \textbf{Interactive Land Area Bounding Box Selector:} Users can click and drag directly across the map canvas to bound any custom agricultural parcel or select preset micro-watershed regions.
    \item \textbf{Visual Overlays:} The recommended pond site is highlighted with an animated, pulsating ripple beacon icon; the contributing catchment drainage basin is rendered as a semi-transparent cyan polygon with dashed boundary styling; and secondary candidate sites are pinned with numbered rank badges.
    \item \textbf{Analytical Engineering Dashboard:} A tabbed sidebar displays real-time KPI metrics, household water security support figures, an interactive monthly rainfall bar chart, dynamic runoff coefficient sliders, and pond cross-section engineering schematics.
\end{itemize}

\subsection{Database, Caching, and Process Storage}
Given the read-heavy nature of spatial contour datasets, the system employs an in-memory caching and pre-computed static asset strategy. The full 6.4~MB contour dataset (2,710 contours) is pre-processed into an optimized 808-feature GeoJSON vector file (\texttt{contours.geojson}) and a baseline hydrological result cache (\texttt{baseline\_result.json}). This enables the initial map load to render in under 50~ms without incurring repetitive 25-second matrix interpolation cycles. Process state and watchdog logs are persistently maintained in \texttt{service.log}.

%% =====================================================================
\section{CSD Themes and Topics Applied in the Project}
\label{sec:csd-themes}
Table~\ref{tab:csd-themes} demonstrates the direct alignment between system design decisions and core Computer System Design (CSD) curriculum topics.

\begin{table*}[t]
  \caption{Mapping of core CSD themes/topics to their concrete implementation in AquaPlan AI}
  \label{tab:csd-themes}
  \begin{tabular}{p{3.2cm}p{5.0cm}p{6.8cm}}
    \toprule
    CSD Theme/Topic & Where used in the project & Justification and Design Rationale \\
    \midrule
    API Design (REST) & \texttt{main.py} (\texttt{/api/analyze-bounds}, \texttt{/api/baseline}) & Stateless, standardized REST contracts decouple client mapping from computational raster pipelines, enabling independent scaling and headless testing. \\
    Load Distribution / Port Mapping & Dual workers on ports 3000 \& 6000 (\texttt{start\_daemon.py}) & Exposes service across IIT Bhilai network at \texttt{http://10.1.75.51:6253} (and \texttt{:3253}), enabling concurrent access across campus subnets. \\
    Caching & Pre-computed \texttt{baseline\_result.json} \& \texttt{contours.geojson} & Eliminates redundant 26-second Delaunay triangulation and D8 routing for standard full-campus queries, reducing query latency by 99.8\%. \\
    Spatial Indexing / Coordinate Cropping & Bounding box spatial filter in \texttt{terrain\_analysis.py} & Filters 2,710 contours down to the user-selected parcel before interpolation, bounding the Delaunay matrix size from $1621 \times 1313$ to $400 \times 400$. \\
    Concurrency / Async Processing & Asynchronous request handling via FastAPI/ASGI & Asynchronous event loop prevents server thread blocking during client map tile streaming while heavy numerical jobs run in worker threads. \\
    Process Supervision \& Daemonization & \texttt{start\_daemon.py} watchdog supervisor & Detects unexpected worker termination, automatically re-spawning Uvicorn processes within 3~seconds to maintain high availability. \\
    Algorithms and Complexity & D8 Flow Accumulation \& Quickhull Convex Hull & Topological sorting of flat elevation arrays enables $O(N \log N)$ flow accumulation, avoiding exponential recursion in cyclic basins. \\
    Modular Design Patterns & Strategy Pattern for runoff models; Layer Pattern for GIS & Isolates coordinate transformation, surface interpolation, hydrology, and presentation into discrete, testable architectural layers. \\
    Fault Tolerance \& Graceful Degradation & Network retry loops in deployment \& fallback sinks & Fallback to largest interior topographical sink if all candidate depressions violate strict river proximity criteria. \\
    Testing Strategy & Automated PowerShell \& Python test harnesses & Benchmarked socket reachability, KML parsing throughput, bounding box subregion latency, and API payload integrity. \\
    Version Control & Git repository initialized in \texttt{pond-planning} & Structured branching, atomic commits, and reproducible package archives ensure configuration traceability. \\
    Network Address Translation (NAT) & IIT Bhilai Port Forwarding (Port 2253 / 6253) & System port mapping formula ($X000 \rightarrow X253$) accommodated to ensure public network reachability on all allotted IP subnets. \\
    Resource Efficiency & Headless vector simplification of contour polylines & Polyline point sampling reduces client memory consumption and prevents WebGL context loss on mobile browsers. \\
    \bottomrule
  \end{tabular}
\end{table*}

\subsubsection*{Engineering Rationale for Significant Architectural Decisions}
Among the CSD principles applied, two design decisions proved especially critical:
\begin{enumerate}
    \item \textbf{Decoupled Pre-computed Baseline Caching vs. Dynamic Sub-grid Interpolation:} Running a full-scale linear Delaunay interpolation over 2,710 contour polylines on a 2-meter grid creates an array of over 2.1 million elements ($1621 \times 1313$). In a synchronous monolith without caching, every initial page visit would stall for $\sim 26.8$~seconds. By decoupling the static full-study-area analysis into a pre-computed JSON baseline while restricting runtime interpolation strictly to user-selected bounding boxes, the system achieves sub-50~ms initial page loads and sub-3~second dynamic parcel analysis.
    \item \textbf{Supervisor Watchdog Daemon with Dual-Port Mapping:} In university and institutional compute clusters, interactive terminal sessions are routinely terminated by SSH inactivity timeouts or PAM session reapers. Implementing \texttt{start\_daemon.py} as a decoupled background supervisor with infinite watchdog monitoring ensures that workers on both port 3000 and port 6000 are continuously revitalized. Furthermore, configuring dual listening ports satisfies institutional port forwarding rules ($10.1.75.51:6253 \rightarrow 6000$ and $10.1.75.51:3253 \rightarrow 3000$), ensuring uninterrupted accessibility for evaluators.
\end{enumerate}

%% =====================================================================
\section{Results and Evaluation}
\label{sec:results}
The platform was thoroughly evaluated on the full topographic dataset of the IIT Bhilai campus (6.4~MB KML, 2,710 contour lines spanning elevation 267.0~m to 298.0~m AMSL). Table~\ref{tab:results-candidates} summarizes the top candidate pond locations identified by the multi-criteria D8 flow accumulation engine.

\begin{table}[h]
  \caption{Evaluated candidate pond sites and hydrological recommendations}
  \label{tab:results-candidates}
  \begin{tabular}{ccccccc}
    \toprule
    Rank & Latitude (°N) & Longitude (°E) & Elevation (m) & Catchment Area ($\text{m}^2$) & Suitability Score & Annual Runoff ($\text{m}^3$) \\
    \midrule
    \textbf{1 (Primary)} & \textbf{21.25205} & \textbf{81.30924} & \textbf{290.0} & \textbf{780.0} & \textbf{0.248} & \textbf{361.6} \\
    2 (Alternative) & 21.25568 & 81.29776 & 286.0 & 776.0 & 0.210 & 359.7 \\
    3 (Alternative) & 21.25028 & 81.30800 & 286.0 & 784.0 & 0.208 & 363.5 \\
    Subregion BBox & 21.25778 & 81.30546 & 293.0 & 664.0 & 0.313 & 307.8 \\
    \bottomrule
  \end{tabular}
\end{table}

\subsection{Performance and Load Benchmarks}
System execution benchmarks recorded across repetitive queries on the allotted system (\texttt{student@10.1.75.51}) are reported below:
\begin{itemize}
    \item \textbf{KML Binary Ingestion:} $0.03\text{ s}$ for 6.4~MB file.
    \item \textbf{Contour Vector Parsing (lxml):} $0.25\text{ s}$ for 2,710 contour features.
    \item \textbf{Full Baseline Query Latency:} $10.2\text{ ms}$ (cached endpoint).
    \item \textbf{Dynamic Sub-grid DEM Interpolation:} $2.54\text{ s}$ for a $963 \times 666$ sub-grid.
    \item \textbf{End-to-End Dynamic Bounding Box API Response:} $7.96\text{ s}$ (including spatial clipping, DEM interpolation, D8 flow routing, flow accumulation, convex hull extraction, and JSON serialization).
    \item \textbf{Memory Footprint:} $\sim 68\text{ MB}$ RSS per Uvicorn worker process.
\end{itemize}

%% =====================================================================
\section{Discussion and Limitations}
\label{sec:discussion}
The experimental results demonstrate that the combination of DEM interpolation and D8 flow routing effectively identifies natural topographic depressions while systematically bypassing high-energy stream channels. The dynamic parameter sliders enable field planners to assess how changes in soil compaction or anomalous drought years impact reservoir recharge. 

However, two operational limitations must be noted: (i) \textbf{Contour Resolution Sensitivity:} Because the input contours have a 1-meter vertical contour interval, micro-topographic variations under 1~meter cannot be fully resolved, occasionally yielding flat saddle zones where nearest-neighbor imputation is required; (ii) \textbf{D8 Single-Direction Heuristic:} In exceptionally flat agricultural fields, the D8 algorithm routes all cell discharge to a single steepest neighbor, whereas dispersive flow models (e.g., D-Infinity or Multiple Flow Direction) could model divergent sheetflow with higher fidelity, albeit at greater computational cost.

%% =====================================================================
\section{AI Tool Usage Declaration}
\label{sec:ai-usage}
In compliance with the assignment's LLM and AI Tool Usage Policy, the authors declare that AI assistance (specifically Google DeepMind Gemini models within the Antigravity IDE) was utilized during the development of this project. AI was employed for: (i) reformatting and styling the ACM \texttt{acmart} LaTeX template; (ii) synthesizing boilerplate REST endpoint schemas in FastAPI; (iii) structuring the Leaflet.js client-side bounding box drawing interaction logic; and (iv) debugging OpenSSH socket connection timeouts and Windows-to-Linux path serialization. All algorithmic implementations of D8 flow routing, mathematical runoff derivations, spatial projections, and scientific evaluations were rigorously reviewed, validated, and tested by the team prior to submission.

%% =====================================================================
\appendix

\section{Source Code, Repository, and Deployment URLs}
\label{sec:appendix}

\subsection{Active Deployment \& Repository URLs}
The system is actively running in the background under supervisor monitoring and can be accessed across the institutional network:
\begin{itemize}
    \item \textbf{GitHub Repository URL:} \url{https://github.com/nikhilesh-git/pond-planning-system}
    \item \textbf{Final Working Front-End URL (Port 6000):} \url{http://10.1.75.51:6253}
    \item \textbf{Backend REST API \& Swagger Docs (Port 3000):} \url{http://10.1.75.51:3253/docs}
    \item \textbf{Health Check Endpoint:} \url{http://10.1.75.51:6253/api/health}
\end{itemize}

\subsection{Repository Folder Structure}
The project codebase is organized within \texttt{/home/student/pond-planning}:
\begin{verbatim}
/home/student/pond-planning/
|-- main.py                     # FastAPI REST API & static web server
|-- terrain_analysis.py         # DEM interpolation, D8 flow routing & hydrology
|-- start_daemon.py             # Watchdog supervisor daemon (ports 3000 & 6000)
|-- start_service.sh            # Service launcher shell script
|-- contours_1m.kml             # Topographic contour map dataset (6.4 MB)
|-- requirements.txt            # Python dependencies (FastAPI, NumPy, SciPy, PyProj)
|-- service.log                 # Continuous supervisor activity and restart log
|-- report.tex                  # ACM manuscript LaTeX final technical report
|-- static/                     # Web GIS Client assets
    |-- index.html              # Modern responsive HTML5 GIS interface
    |-- style.css               # Glassmorphic dark design system & styling
    |-- app.js                  # Leaflet map controller & client-side hydrology
    |-- contours.geojson        # Vector contour lines for overlay visualization
    |-- baseline_result.json    # Cached full study area baseline analysis
    `-- bounds.json             # Geographic bounds metadata of study area
\end{verbatim}

\end{document}
\endinput
